\section{Scope and Sharpness of the Hypotheses}
\label{sec:scope-and-sharpness}

The hypotheses in \eqref{eq:ng-iid-assumptions} include centered,
variance-normalized Laplace densities, despite their cusp at the origin.
They also include variance-normalized densities proportional to $e^{-x^4}$.
Symmetry is not needed in any step of the proof.

For a family allowing both asymmetry and heavy tails, take
\begin{equation}
 \widetilde f_{\nu,\beta}(y)
   =c_{\nu,\beta}(1+y^2)^{-(\nu+1)/2}
                           e^{\beta\arctan y},
 \qquad \nu>2,\quad \beta\in\mathbb R,
 \label{eq:ng-asymmetric-heavy-tail-example}
\end{equation}
and center and rescale to variance one.
The unstandardized location score is
\[
 \widetilde s(y)=\frac{(\nu+1)y-\beta}{1+y^2}.
\]
Both \(\widetilde s\) and \(y\widetilde s\) are bounded.
The density is positive and smooth.
It has a finite second moment.
Centering and rescaling preserve finiteness of the location and scale
information.
So the standardized density satisfies \eqref{eq:ng-iid-assumptions}.
For \(2<\nu\le4\), its fourth moment is infinite.
When \(\beta\ne0\), the two tails have unequal leading constants.
So centering does not make the density even.
The case \(\beta=0\) is a rescaled Student family.
These examples show that neither evenness nor a fourth moment is implicit in
the theorem.

The Gamma examples used to illustrate the Fourier response belong to
\cref{lem:ng-fourier-ode}.
Their centered densities still vanish on a half-line.
So they do not satisfy the full-support hypothesis of
\cref{thm:ng-iid-endpoint}.
A nonzero Fourier response establishes local detection.
Alone, this does not establish the general measure-equivalence theorem for a
density with restricted support.

The second moment has a specific role beyond differentiating one Fourier
feature.
It makes every \(\ell^2\) row into a convergent random series.
It supplies uniform integrability of squares of linear forms.
It gives the Gaussian-residual compactness argument.
So the first-moment ODE result is stronger in its moment scope than the
arbitrary-operator equivalence theorem.
Extending the latter to infinite-variance coordinates requires a different
domain of admissible rows and a corresponding operator criterion.
For example, for symmetric nondegenerate \(\alpha\)-stable coordinates,
\(0<\alpha<2\), convergence of a row series requires
$\sum_j|a_j|^\alpha<\infty$: its partial-sum characteristic
functions are
\[
 \exp\left(-c|t|^\alpha\sum_{j\le n}|a_j|^\alpha\right).
\]
An arbitrary \(\ell^2\) row need not meet that requirement.
The condition on admissible rows must also be satisfied when using bounded
or fractional Fourier features.

The Hellinger path estimates give $L^1$ convergence and uniform
integrability of likelihoods.
The finite-frequency functions are bounded in each coordinate.
Their dimension-independent bounds are the $L^2(\mu_n)$ and generator
estimates proved above.
No variance bound for their aggregate under an arbitrary mixed output law is
needed in the tail-symmetry argument.
